%% ============================================================
%% Appendix: Paper Quality Statistics
%% ============================================================

\section{Paper Quality Statistics}
\label{app:paper_quality}

We report structural statistics of all 75~papers (5~systems $\times$ 3~seeds $\times$ 5~tasks) to characterize the surface-level quality of AI-generated manuscripts.
Statistics are extracted automatically from the LaTeX source and compiled PDF of each paper.


\Cref{tab:paper_quality_summary} reports the mean and standard deviation of each metric across the 15~papers per system.
ARC papers are the longest (12.7~$\pm$~1.2 pages, 5{,}417~$\pm$~863 words) and most figure-heavy (4.3~$\pm$~2.0 per paper).
\sys{} has the largest bibliography pool (55.3~$\pm$~3.6 entries) and 3.8~$\pm$~0.4 figures per paper, though only 18.3~$\pm$~4.5 keys are actually cited---fewer than AIR's 19.5~$\pm$~3.9 unique citation keys.
AIR is equation-heavy (7.1~$\pm$~1.8 per paper) and table-heavy (4.9~$\pm$~1.5) but generates nearly zero figures (0.1~$\pm$~0.2).
Sakana produces the shortest papers (5.8~$\pm$~0.8 pages, 2{,}606~$\pm$~490 words) with zero figures across all 15~papers; DS is similarly brief (6.7~$\pm$~2.1 pages).

\begin{table}[h]
\centering
\small
\caption{Paper quality statistics across five systems (mean $\pm$ std over 15~papers each).}
\label{tab:paper_quality_summary}
\begin{tabular}{@{}lccccc@{}}
\toprule
\textbf{Metric} & \textbf{ARC} & \textbf{DS} & \textbf{AIR} & \textbf{Sakana} & \textbf{\sys{}} \\
\midrule
Pages & 12.7 $\pm$ 1.2 & 6.7 $\pm$ 2.1 & 10.4 $\pm$ 1.1 & 5.8 $\pm$ 0.8 & 10.2 $\pm$ 0.8 \\
Words & 5{,}417 $\pm$ 863 & 2{,}741 $\pm$ 968 & 5{,}116 $\pm$ 524 & 2{,}606 $\pm$ 490 & 4{,}643 $\pm$ 438 \\
Figures & 4.3 $\pm$ 2.0 & 1.3 $\pm$ 0.4 & 0.1 $\pm$ 0.2 & 0.0 $\pm$ 0.0 & 3.8 $\pm$ 0.4 \\
Tables & 2.1 $\pm$ 0.7 & 0.8 $\pm$ 0.7 & 4.9 $\pm$ 1.5 & 2.3 $\pm$ 1.1 & 1.9 $\pm$ 0.2 \\
Equations & 1.3 $\pm$ 1.6 & 1.9 $\pm$ 2.0 & 7.1 $\pm$ 1.8 & 2.2 $\pm$ 1.5 & 3.0 $\pm$ 1.1 \\
Unique cite keys & 23.5 $\pm$ 8.7 & 9.3 $\pm$ 6.9 & 19.5 $\pm$ 3.9 & 10.5 $\pm$ 2.5 & 18.3 $\pm$ 4.5 \\
Bib entries & 23.5 $\pm$ 8.7 & 13.1 $\pm$ 13.2 & 15.8 $\pm$ 1.9 & 10.6 $\pm$ 2.6 & 55.3 $\pm$ 3.6 \\
Sections & 6.6 $\pm$ 2.5 & 6.5 $\pm$ 1.0 & 5.0 $\pm$ 0.0 & 6.5 $\pm$ 0.7 & 5.5 $\pm$ 0.6 \\
Subsections & 9.1 $\pm$ 3.9 & 9.7 $\pm$ 3.4 & 9.9 $\pm$ 1.7 & 6.0 $\pm$ 1.1 & 7.5 $\pm$ 1.4 \\
\bottomrule
\end{tabular}
\end{table}

\subsection{Failure Mode Case Studies}
\label{sec:exp:failures}

We highlight four failure modes from the 75-paper audit, each illustrating a different evidence chain break that \coeaudit{}'s integrity checks are designed to catch.

\paragraph{Case 1: Six orders of magnitude (ARC, LLM-SQL, seed~2).}
The paper introduces ``SCOR,'' a static column-ordering routine, and reports a combined score of 1{,}538{,}006.69---on a benchmark whose scoring metric uses a $[0, 1]$ scale.
The number is not a transcription error: it is the sum of squared prefix-hit lengths across datasets, an internal metric that the system computed and presented as if it were the ADRS score.
The paper is internally coherent---it defines its own evaluation protocol, runs controlled comparisons against a baseline (scoring 1{,}537{,}927.99), and draws reasonable conclusions within that framing.
An automated reviewer evaluating narrative quality alone would find nothing wrong.
Score Verification catches this immediately: the canonical evaluator re-run crashes (the submitted code fails to produce a valid solution), making the entire evidence chain from score to evaluator unresolvable.

\paragraph{Case 2: A bibliography from model memory (AIR, PRISM, seed~1).}
This paper's bibliography contains 15~references.
Reference Verification finds that 3 of those references are hallucinated---we could not find matching publications in Semantic Scholar, arXiv, or other scholarly databases.
This is not an edge case---AIR and DS produce hallucinated references at rates of 9\% and 21\% respectively (Table~\ref{tab:integrity}), compared to 0\% for systems with structured retrieval pipelines.

\paragraph{Case 3: Convergent specification violation (DS, LLM-SQL, seed~1).}
Under the I2 majority-vote protocol ($K{=}5$), 2 of 5 judges flag this paper---below the majority threshold, so it is not counted as a violation in Table~\ref{tab:integrity}.
Nevertheless, the submitted code achieves a legitimate score of 0.697 that passes Score Verification, but it does so by exploiting a gap between what the evaluator checks and what the benchmark intends to measure.
The code sorts columns differently per row-group block, then renames all columns back to the original schema before concatenation---this makes \texttt{pd.concat} assemble blocks in insertion order rather than realigning by column name, effectively permuting column order per row-group.
The evaluator validates row counts and total character counts but not column-to-column correspondence, so the permutation goes undetected.
The same exploit appears independently in two other systems (AIR seed~1 and \sys{} seed~2), providing convergent evidence that this is a genuine benchmark vulnerability rather than an isolated accident.%
\footnote{For \sys{} seed~2, the exploit is present in the submitted code, but the I2 majority-vote protocol did not reach consensus (1 of 5 judges flagged it), so it is not counted as a violation in Table~\ref{tab:integrity}. This illustrates a limitation of LLM-judged checks at the current vote threshold.}

\paragraph{Case 4: Near-correct score, fictional algorithm (ARC, TXN, seed~1).}
This paper nearly passes Score Verification: the reported score of 3{,}311 is within 3\% of the canonical evaluator re-run mean (3{,}214), just outside the adaptive tolerance threshold.
Yet Method-Code Alignment reveals a complete disconnect between what the paper describes and what the code implements.
The paper introduces ``STAR,'' a system built on bitwise integer encoding for conflict detection, an $O(1)$ surrogate cost model, and equidistant placement of high-contention anchor transactions.
The submitted code implements none of these: it uses standard Python sets for conflict tracking, calls the full simulator on every iteration (no surrogate), and clusters read-heavy keys sequentially rather than distributing write-heavy anchors.
This case illustrates why Score Verification alone is insufficient---the solver works, but the paper describes a different algorithm entirely, making the method section unreproducible regardless of how accurate the reported numbers are.


\begin{comment}
\subsection{Per-Paper Details}

\Cref{tab:paper_quality_detail} lists statistics for each of the 75~papers.
Task abbreviations: CC = Cloudcast, EPLB = Expert-Parallel Load Balancing, SQL = LLM-SQL, TXN = Transaction Scheduling.

\begin{longtable}{@{}lll rrrr rrrr@{}}
\caption{Per-paper statistics for all 75~papers. Size is in KB.}
\label{tab:paper_quality_detail} \\
\toprule
\textbf{System} & \textbf{Sd} & \textbf{Task} & \textbf{Pg} & \textbf{Words} & \textbf{KB} & \textbf{Fig} & \textbf{Tbl} & \textbf{Eq} & \textbf{Cites} & \textbf{Bib} \\
\midrule
\endfirsthead
\multicolumn{11}{@{}l}{\small\itshape\tablename~\thetable{}, continued} \\
\toprule
\textbf{System} & \textbf{Sd} & \textbf{Task} & \textbf{Pg} & \textbf{Words} & \textbf{KB} & \textbf{Fig} & \textbf{Tbl} & \textbf{Eq} & \textbf{Cites} & \textbf{Bib} \\
\midrule
\endhead
\midrule
\multicolumn{11}{r@{}}{\small\itshape Continued on next page} \\
\endfoot
\bottomrule
\endlastfoot
%% Sakana
Sakana & 1 & CC    & 6 & 2{,}618 & 143 & 0 & 1 & 3 & 8 & 8 \\
Sakana & 1 & EPLB  & 6 & 2{,}348 & 140 & 0 & 4 & 1 & 13 & 13 \\
Sakana & 1 & SQL   & 5 & 2{,}239 & 148 & 0 & 2 & 2 & 9 & 9 \\
Sakana & 1 & Prism & 8 & 4{,}076 & 170 & 0 & 4 & 4 & 11 & 11 \\
Sakana & 1 & TXN   & 6 & 2{,}918 & 151 & 0 & 2 & 2 & 14 & 14 \\
Sakana & 2 & CC    & 5 & 2{,}225 & 133 & 0 & 2 & 1 & 8 & 8 \\
Sakana & 2 & EPLB  & 5 & 2{,}204 & 125 & 0 & 4 & 1 & 14 & 14 \\
Sakana & 2 & SQL   & 5 & 2{,}369 & 170 & 0 & 2 & 3 & 9 & 9 \\
Sakana & 2 & Prism & 7 & 3{,}010 & 149 & 0 & 2 & 6 & 10 & 10 \\
Sakana & 2 & TXN   & 6 & 2{,}427 & 120 & 0 & 4 & 1 & 13 & 13 \\
Sakana & 3 & CC    & 6 & 2{,}933 & 137 & 0 & 2 & 1 & 6 & 6 \\
Sakana & 3 & EPLB  & 6 & 2{,}547 & 137 & 0 & 2 & 2 & 11 & 12 \\
Sakana & 3 & SQL   & 5 & 2{,}056 & 146 & 0 & 2 & 3 & 7 & 7 \\
Sakana & 3 & Prism & 6 & 2{,}853 & 169 & 0 & 1 & 3 & 11 & 12 \\
Sakana & 3 & TXN   & 5 & 2{,}273 & 141 & 0 & 1 & 0 & 13 & 13 \\
\midrule
%% ARC
ARC & 1 & CC    & 13 & 5{,}828 & 461 & 2 & 2 & 1 & 24 & 24 \\
ARC & 1 & EPLB  & 12 & 4{,}692 & 1{,}922 & 6 & 2 & 0 & 20 & 20 \\
ARC & 1 & SQL   & 9 & 3{,}448 & 460 & 4 & 2 & 0 & 0 & 0 \\
ARC & 1 & Prism & 14 & 6{,}875 & 389 & 1 & 3 & 0 & 31 & 31 \\
ARC & 1 & TXN   & 13 & 6{,}455 & 400 & 1 & 3 & 5 & 35 & 35 \\
ARC & 2 & CC    & 14 & 5{,}997 & 745 & 5 & 2 & 3 & 24 & 24 \\
ARC & 2 & EPLB  & 13 & 5{,}606 & 1{,}086 & 7 & 1 & 0 & 21 & 21 \\
ARC & 2 & SQL   & 12 & 4{,}426 & 1{,}623 & 7 & 2 & 2 & 21 & 21 \\
ARC & 2 & Prism & 13 & 6{,}476 & 430 & 2 & 2 & 1 & 33 & 33 \\
ARC & 2 & TXN   & 14 & 5{,}424 & 971 & 6 & 3 & 2 & 35 & 35 \\
ARC & 3 & CC    & 12 & 5{,}041 & 731 & 4 & 2 & 0 & 17 & 17 \\
ARC & 3 & EPLB  & 13 & 5{,}275 & 681 & 7 & 1 & 0 & 20 & 20 \\
ARC & 3 & SQL   & 12 & 4{,}594 & 565 & 4 & 1 & 2 & 16 & 16 \\
ARC & 3 & Prism & 13 & 5{,}710 & 905 & 4 & 3 & 4 & 26 & 26 \\
ARC & 3 & TXN   & 14 & 5{,}412 & 586 & 4 & 2 & 0 & 29 & 29 \\
\midrule
%% DS
DS & 1 & CC    & 5 & 1{,}781 & 149 & 1 & 1 & 1 & 21 & 32 \\
DS & 1 & EPLB  & 8 & 3{,}741 & 136 & 1 & 1 & 0 & 23 & 41 \\
DS & 1 & SQL   & 5 & 2{,}213 & 252 & 1 & 1 & 0 & 0 & 3 \\
DS & 1 & Prism & 8 & 3{,}348 & 138 & 1 & 2 & 1 & 16 & 32 \\
DS & 1 & TXN   & 9 & 4{,}250 & 154 & 2 & 2 & 1 & 10 & 9 \\
DS & 2 & CC    & 8 & 3{,}960 & 173 & 1 & 1 & 4 & 1 & 1 \\
DS & 2 & EPLB  & 8 & 2{,}412 & 390 & 1 & 0 & 1 & 13 & 30 \\
DS & 2 & SQL   & 4 & 1{,}197 & 94 & 1 & 0 & 1 & 5 & 5 \\
DS & 2 & Prism & 5 & 1{,}968 & 162 & 2 & 1 & 5 & 8 & 8 \\
DS & 2 & TXN   & 7 & 2{,}950 & 164 & 2 & 0 & 3 & 6 & 6 \\
DS & 3 & CC    & 6 & 2{,}716 & 146 & 1 & 1 & 1 & 5 & 5 \\
DS & 3 & EPLB  & 12 & 4{,}078 & 224 & 1 & 1 & 2 & 17 & 17 \\
DS & 3 & SQL   & 5 & 1{,}873 & 121 & 1 & 0 & 0 & 3 & 3 \\
DS & 3 & Prism & 7 & 3{,}234 & 179 & 2 & 0 & 7 & 5 & 5 \\
DS & 3 & TXN   & 4 & 1{,}398 & 130 & 1 & 0 & 1 & 7 & 0 \\
\midrule
%% AIR
AIR & 1 & CC    & 9 & 4{,}596 & 181 & 0 & 3 & 9 & 18 & 14 \\
AIR & 1 & EPLB  & 13 & 6{,}262 & 216 & 1 & 7 & 7 & 18 & 16 \\
AIR & 1 & SQL   & 9 & 4{,}760 & 182 & 0 & 3 & 5 & 17 & 15 \\
AIR & 1 & Prism & 10 & 5{,}031 & 188 & 0 & 4 & 5 & 17 & 15 \\
AIR & 1 & TXN   & 10 & 4{,}899 & 183 & 0 & 4 & 5 & 30 & 22 \\
AIR & 2 & CC    & 11 & 5{,}413 & 196 & 0 & 8 & 8 & 18 & 15 \\
AIR & 2 & EPLB  & 10 & 4{,}810 & 210 & 0 & 3 & 5 & 15 & 15 \\
AIR & 2 & SQL   & 11 & 4{,}905 & 204 & 0 & 6 & 7 & 18 & 15 \\
AIR & 2 & Prism & 10 & 5{,}222 & 201 & 0 & 4 & 10 & 17 & 14 \\
AIR & 2 & TXN   & 11 & 5{,}360 & 211 & 0 & 5 & 7 & 20 & 15 \\
AIR & 3 & CC    & 9 & 4{,}662 & 196 & 0 & 4 & 8 & 17 & 16 \\
AIR & 3 & EPLB  & 9 & 4{,}332 & 158 & 0 & 5 & 4 & 23 & 17 \\
AIR & 3 & SQL   & 12 & 6{,}050 & 204 & 0 & 6 & 8 & 24 & 16 \\
AIR & 3 & Prism & 11 & 4{,}805 & 201 & 0 & 7 & 10 & 24 & 17 \\
AIR & 3 & TXN   & 11 & 5{,}632 & 181 & 0 & 5 & 8 & 17 & 15 \\
\midrule
%% ScientistOne
\sys{} & 1 & CC    & 11 & 5{,}175 & 2{,}865 & 4 & 2 & 2 & 21 & 56 \\
\sys{} & 1 & EPLB  & 9 & 3{,}631 & 2{,}378 & 3 & 1 & 4 & 10 & 52 \\
\sys{} & 1 & SQL   & 9 & 4{,}196 & 2{,}124 & 3 & 2 & 4 & 21 & 62 \\
\sys{} & 1 & Prism & 11 & 5{,}263 & 2{,}716 & 4 & 2 & 3 & 27 & 55 \\
\sys{} & 1 & TXN   & 10 & 4{,}714 & 2{,}961 & 4 & 2 & 4 & 14 & 54 \\
\sys{} & 2 & CC    & 11 & 4{,}877 & 3{,}263 & 4 & 2 & 2 & 18 & 53 \\
\sys{} & 2 & EPLB  & 11 & 4{,}794 & 3{,}140 & 4 & 2 & 3 & 20 & 54 \\
\sys{} & 2 & SQL   & 11 & 5{,}007 & 3{,}209 & 4 & 2 & 1 & 21 & 61 \\
\sys{} & 2 & Prism & 11 & 5{,}061 & 3{,}560 & 4 & 2 & 4 & 24 & 54 \\
\sys{} & 2 & TXN   & 11 & 4{,}870 & 2{,}828 & 4 & 2 & 3 & 14 & 53 \\
\sys{} & 3 & CC    & 10 & 4{,}467 & 2{,}731 & 4 & 2 & 1 & 18 & 54 \\
\sys{} & 3 & EPLB  & 9 & 4{,}116 & 2{,}885 & 3 & 2 & 2 & 17 & 51 \\
\sys{} & 3 & SQL   & 10 & 4{,}666 & 2{,}540 & 4 & 2 & 4 & 22 & 63 \\
\sys{} & 3 & Prism & 10 & 4{,}641 & 2{,}631 & 4 & 2 & 4 & 16 & 55 \\
\sys{} & 3 & TXN   & 9 & 4{,}162 & 3{,}743 & 4 & 2 & 4 & 12 & 53 \\
\end{longtable}
\end{comment}